\section{Results}
\label{sec:results}

% Aggregate success motivates the diagnosis but does not explain the failure mechanism.
We organize the results around four claims: (P1) the failure is concentrated \emph{after} contact, not at reaching; (P2) this post-contact failure is shared across URL methods; (P3) it is not explained by dataset support; and (P4) its cause is interaction-mode aliasing rather than poor geometric resolution.

Table~\ref{tab:overall_state} reports aggregate state-based success. \gciql{} leads at $68.3\%$ IQM, \fb{} reaches $50.0\%$, and the contrastive and predictive representations trail at $25.9\%$ (\crl{}) and $22.7\%$ (\rldp{}). This ordering motivates the diagnosis but does not, by itself, identify whether a method fails before contact, during contact acquisition, or after the cube is lifted.

\begin{table}[htb]
\centering
\caption{Aggregate state-based evaluation success. Values are IQM $\pm$ half-width of the $95\%$ stratified-bootstrap confidence interval across $10$ seeds.\protect\footnotemark}
\label{tab:overall_state}
\begin{tabular}{lcccccc}
\toprule
\textbf{Method} & \textbf{task1} & \textbf{task2} & \textbf{task3} & \textbf{task4} & \textbf{task5} & \textbf{Overall} \\
\midrule
\fb{}    & $16.7 \pm 10.0$ & $68.3 \pm 10.9$ & $75.0 \pm 10.0$ & $40.0 \pm 13.3$ & $43.3 \pm 18.4$ & $\mathbf{50.0 \pm 7.0}$ \\
\gciql{} & $73.7 \pm 5.9$  & $70.0 \pm 7.0$  & $73.0 \pm 7.0$  & $58.0 \pm 5.2$  & $64.3 \pm 7.0$  & $\mathbf{68.3 \pm 3.7}$ \\
\crl{}   & $24.0 \pm 3.3$  & $28.0 \pm 6.7$  & $28.0 \pm 4.8$  & $24.3 \pm 3.5$  & $26.3 \pm 5.8$  & $\mathbf{25.9 \pm 2.0}$ \\
\rldp{}  & $28.3 \pm 10.8$ & $16.7 \pm 6.7$  & $35.0 \pm 15.9$ & $13.3 \pm 12.5$ & $18.3 \pm 10.0$ & $\mathbf{22.7 \pm 6.0}$ \\
\bottomrule
\end{tabular}
\end{table}
\footnotetext{Aggregate success uses the stochastic evaluation policy (IQM over seeds). The phase decomposition below uses deterministic closed-loop rollouts, so per-method success differs slightly (e.g. \fb{} $49.4\%$, \gciql{} $69.2\%$); we report each metric consistently within its own analysis.}

\subsection{Failures occur after contact, not before it (P1)}

% Reach-before-drop: maintain/transport failures are loss of control before the goal.
Our phase labels operationalize the reach-before-drop event of Section~\ref{sec:theory}: an approach or grasp failure means the agent never establishes control, whereas a maintain-grasp or transport failure means the agent established control and then lost it (entered $D$) before reaching the goal $G_g$. Figure~\ref{fig:funnel} traces this funnel. Every method reaches the cube in essentially all episodes and secures a grasp in $76$--$90\%$, but success then collapses. The separation between methods opens \emph{between} the secured and success stages --- that is, after contact has been established, not before it.

\begin{figure}[htb]
\centering
\includegraphics[width=0.62\linewidth]{f1_funnel.png}
\caption{Phase funnel across methods (deterministic rollouts, $10$ seeds). All methods reach the cube and most secure a grasp; success collapses afterwards. The between-method gap opens at the secured-to-success transition, i.e. after contact.}
\label{fig:funnel}
\end{figure}

\subsection{The post-contact failure is shared across URL methods (P2)}

% Cross-method decomposition: the gap is in maintain/transport, not approach.
Table~\ref{tab:cross_method_template} and Figure~\ref{fig:composition} give the terminal-outcome distribution for each method. Approach failure is negligible everywhere ($\leq 0.4\%$); the methods differ only in how their failures split across grasp, maintain-grasp, and transport. Relative to \gciql{}, the \fb{} success gap of $19.8\pp$ decomposes as
\begin{equation}
    -19.8 = \underbrace{+0.2}_{\text{approach}} \; \underbrace{+2.4}_{\text{grasp}} \; \underbrace{-14.6}_{\text{maintain}} \; \underbrace{-7.8}_{\text{transport}},
\end{equation}
i.e. it is entirely a post-contact deficit; \fb{} in fact has \emph{fewer} approach and grasp failures than \gciql{}. A near-miss check confirms the interpretation: across the \fb{} and \gciql{} rollouts there are no episodes that finish holding the cube within one $5$~cm success radius of the goal, so transport failures are not fine placement errors but failures to bring the held object to the target.

\begin{table}[htb]
\centering
\caption{Cross-method terminal-outcome distribution (S0, \% of episodes, mean across $10$ seeds; rows sum to $100\%$). \fb{}, \gciql{}, \crl{}, and \rldp{} are completed measurements; \psm{} and \dreamer{} indicate planned diagnostics under the same protocol.}
\label{tab:cross_method_template}
\begin{tabular}{lccccc}
\toprule
\textbf{Method} & \textbf{Success} & \textbf{Approach fail} & \textbf{Grasp fail} & \textbf{Maintain fail} & \textbf{Transport fail} \\
\midrule
\fb{}    & $49.4$ & $0.0$ & $9.8$  & $23.8$ & $17.0$ \\
\gciql{} & $69.2$ & $0.2$ & $12.2$ & $9.2$  & $9.2$ \\
\crl{}   & $23.2$ & $0.4$ & $23.8$ & $28.6$ & $24.0$ \\
\rldp{}  & $25.4$ & $0.2$ & $9.6$  & $44.0$ & $20.8$ \\
\psm{}   & \todo{basis coverage} & \todo{basis coverage} & \todo{basis coverage} & \todo{basis coverage} & \todo{basis coverage} \\
\dreamer{} & \todo{model error} & \todo{model error} & \todo{model error} & \todo{model error} & \todo{model error} \\
\bottomrule
\end{tabular}
\end{table}

\begin{figure}[htb]
\centering
\includegraphics[width=0.66\linewidth]{f1_composition.png}
\caption{Outcome composition per method (visual form of Table~\ref{tab:cross_method_template}). The failure mass sits in maintain-grasp and transport for every method.}
\label{fig:composition}
\end{figure}

The pattern is not specific to \fb{}. \rldp{} secures the cube most reliably of any method (lowest grasp failure, $9.6\%$) yet drops it most often (maintain failure $44.0\%$): nearly half of its episodes lift the cube and then lose it. \crl{} fails more uniformly across grasp, maintain, and transport. In all three lower-scoring methods the deficit relative to \gciql{} is dominated by post-contact outcomes. Because the phase labels never read a method's value, representation, or model, the same protocol extends unchanged to \psm{} (test whether any proto-policy covers stable grasped transport) and to model-based methods (test whether the model predicts lift, slip, and drop transitions).

The next two diagnostics group individual states rather than episodes. Because both the play data and the policy rollouts cycle through pick-and-place many times per trajectory, we pool states into three instantaneous regimes defined by gripper and lift state: \emph{approach} (pre-contact), \emph{pickup} (gripper closed, cube on the table), and \emph{carry} (cube lifted and held). These regimes are distinct from the terminal failure phases of Table~\ref{tab:cross_method_template}: a maintain-grasp failure (cube dropped) and a transport failure (cube held but never delivered) are separated only by how an episode ends, and both originate within the carry regime.

\subsection{Dataset support does not explain the failure (P3)}

% Failed states are not uniformly farther from the offline data.
A natural confound is that post-contact failures are simply out-of-distribution. Table~\ref{tab:state_coverage} reports the nearest-neighbor distance $m_{\mathrm{dist}}$ from rollout states to comparable offline states, by regime and outcome, for all four methods. If failure were a pure coverage effect, failed states should be markedly farther from the data than successful ones in every regime. Instead the gap is small in the approach regime and concentrated in the contact regimes (pickup, carry): failed states are modestly farther for \fb{} and \rldp{}, and more so for \gciql{} and \crl{}. The gap appears even for \gciql{}, which mostly succeeds, so it tracks the under-representation of contact transitions in the play data rather than the failure itself, and it is far too small to account for the $20$--$45\pp$ success deficits in Table~\ref{tab:cross_method_template}. Coverage is therefore at most a contributing factor, not the mechanism.

\begin{table}[htb]
\centering
\caption{Dataset-manifold support. $m_{\mathrm{dist}}$ is the nearest-neighbor distance from rollout states to comparable offline states (mean across seeds; larger means lower support). S = success, Tf = transport-fail episodes; regimes are the per-state instantaneous regimes (Section~\ref{sec:results}).}
\label{tab:state_coverage}
\begin{tabular}{lcccccc}
\toprule
 & \multicolumn{2}{c}{Approach} & \multicolumn{2}{c}{Pickup} & \multicolumn{2}{c}{Carry} \\
\cmidrule(lr){2-3}\cmidrule(lr){4-5}\cmidrule(lr){6-7}
\textbf{Method} & \textbf{S} & \textbf{Tf} & \textbf{S} & \textbf{Tf} & \textbf{S} & \textbf{Tf} \\
\midrule
\fb{}    & $2.23$ & $2.28$ & $2.33$ & $2.69$ & $2.09$ & $2.27$ \\
\gciql{} & $2.56$ & $2.85$ & $2.67$ & $3.37$ & $2.64$ & $3.20$ \\
\crl{}   & $2.45$ & $3.02$ & $2.43$ & $3.82$ & $2.54$ & $3.22$ \\
\rldp{}  & $2.30$ & $2.32$ & $2.38$ & $2.99$ & $2.23$ & $2.66$ \\
\bottomrule
\end{tabular}
\end{table}

\subsection{The bottleneck is interaction-mode aliasing, not geometry (P4)}

% Geometry resolution is not the bottleneck: RLDP resolves geometry yet drops more.
Section~\ref{sec:theory} predicts that the failure is an \emph{interaction-mode} gap: observation-close states with different controllability are aliased by an observation-smooth approximator, including the low-rank inner-product $F^\top B$. A natural first guess is that the representation simply fails to resolve where the cube is. Table~\ref{tab:geometry_control} shows this is not the explanation. Probing the backward embedding $B$ for the cube-to-goal distance over carry-regime states, \rldp{}'s $B$ encodes the terminal geometry far better than \fb{}'s (nonlinear $R^2$ $0.75$ vs $0.21$) --- yet \rldp{} fails to \emph{maintain} the grasp nearly twice as often ($44.0\%$ vs $23.8\%$). Geometry in the representation does not translate into control.

\begin{table}[htb]
\centering
\caption{Carry-regime resolution of the backward embedding $B$ for cube-to-goal distance $d$, under a linear and a nonlinear (MLP) readout, with the maintain-grasp failure rate. Better geometric resolution (\rldp{}) coincides with \emph{worse} grasp maintenance.}
\label{tab:geometry_control}
\begin{tabular}{lccc}
\toprule
\textbf{Method} & \textbf{$R^2(B\to d)$ linear} & \textbf{$R^2(B\to d)$ MLP} & \textbf{Maintain-grasp fail (\%)} \\
\midrule
\fb{}   & $0.20$ & $0.21$ & $23.8$ \\
\rldp{} & $0.55$ & $\mathbf{0.75}$ & $44.0$ \\
\bottomrule
\end{tabular}
\end{table}

What the learned value does fail to provide is a transport signal. Figure~\ref{fig:value_vs_dist} plots the extracted value against cube-to-goal distance by phase: a useful critic would rise as the cube approaches the goal, but in the carry regime value is flat or non-monotone for every method. Consistent with the proposition, even when the geometry is present in $B$ (\rldp{}), the bilinear value that reads $B$ linearly cannot exploit it: \rldp{}'s $B$ is linear-$R^2$ $0.55$ but nonlinear-$R^2$ $0.75$, so the inner-product critic is the binding constraint. The diagnosis is therefore not that the cube is hard to locate, but that observation-smooth, low-rank value approximators alias the controllable and about-to-fail contact states that the policy must distinguish to preserve and exploit the grasp. Table~\ref{tab:aliasing} makes the aliasing direct. For each fail-bound state we find its nearest success-bound state \emph{in observation space} and ask whether the value ranks the success-bound one higher; a value that resolved controllability would score near $1$. Every method stays near chance ($0.48$--$0.60$) in every regime, even though the matched states are genuine observation neighbors. The value cannot separate look-alike states that will succeed from those that will fail --- exactly the interaction-mode aliasing the proposition predicts.

\begin{figure}[htb]
\centering
\includegraphics[width=\linewidth]{f3_value_vs_dist.png}
\caption{Extracted value versus cube-to-goal distance on training states, by phase ($z$-scored per method). A transport signal would rise as distance $\to 0$; instead value is flat or non-monotone in the carry regime for all methods.}
\label{fig:value_vs_dist}
\end{figure}

\begin{table}[htb]
\centering
\caption{Matched-observation aliasing test: for each fail-bound state, the fraction of nearest success-bound states (in observation space) that the value ranks higher, by regime, across seeds. $0.5$ = the value cannot separate observation-close controllable vs about-to-fail states. Matched pairs lie well within the data spread (standardized distance ${\approx}2.7$ vs ${\approx}7$ median).}
\label{tab:aliasing}
\begin{tabular}{lccc}
\toprule
\textbf{Method} & \textbf{Approach} & \textbf{Pickup} & \textbf{Carry} \\
\midrule
\fb{}    & $0.59$ & $0.57$ & $0.59$ \\
\gciql{} & $0.50$ & $0.48$ & $0.50$ \\
\crl{}   & $0.55$ & $0.55$ & $0.55$ \\
\rldp{}  & $0.59$ & $0.57$ & $0.60$ \\
\bottomrule
\end{tabular}
\end{table}
